\section{Modelo relacional y documental}
\subsection{Estructura relacional normalizada}
El esquema de entrega se define en
\archivo{sprint1/postgres/schema_v2_3nf.sql} \cite{pg3nf}. Conserva las doce entidades del
núcleo anterior y añade dos relaciones auxiliares: tipos de dispositivo y
asignaciones sensor--canal. El modelo definitivo contiene catorce tablas.
El esquema de práctica \archivo{sprint1}, que usa \archivo{readings} y
\archivo{storage_id}, es independiente y no representa este modelo.

\begin{longtable}{>{\raggedright\arraybackslash}p{0.25\textwidth}>{\raggedright\arraybackslash}p{0.26\textwidth}>{\raggedright\arraybackslash}p{0.40\textwidth}}
\toprule Tabla & Claves candidatas & Función \\
\midrule\endhead
sources & source\_id & Dataset y publicador. \\
source\_versions & source\_version & Snapshot y fuente. \\
zones & zone\_id & Zona administrativa y distrito padre opcional. \\
locations & location\_id & Nombre del sitio, coordenadas y zona opcional. \\
stations & station\_id; fuente/código & Sitio lógico del inventario. \\
sensor\_categories & category\_id & Clasificación administrativa de dispositivos. \\
measurement\_types & measurement\_type; fuente/código & Magnitud medida y unidad; independiente del tipo físico. \\
organizations & organization\_id & Organismo y responsable. \\
users & user\_id; email & Usuario de aplicación. \\
sensor\_types & sensor\_type & Tipo de dispositivo y categoría. \\
sensors & sensor\_id & Dispositivo administrativo, instalación y estado. \\
channels & channel\_id; estación/clave & Serie lógica de una magnitud. \\
channel\_sensors & channel\_id & Asignación conocida de canal a dispositivo. \\
observations & measurement\_id; canal/UTC & Hecho observado y procedencia variable. \\
\bottomrule
\end{longtable}

\clearpage
\begin{figure}[p]
\centering
\begin{tikzpicture}[x=3.65cm,y=2.1cm,
 entidad/.style={draw,fill=white,rounded corners,align=center,font=\scriptsize,minimum width=2.9cm,minimum height=0.75cm},
 relacion/.style={-{Latex},thin}]
\node[entidad] (src) at (0,0) {sources};
\node[entidad] (ver) at (1,0) {source\_versions};
\node[entidad] (org) at (2,0) {organizations};
\node[entidad] (usr) at (3,0) {users};
\node[entidad] (sta) at (0,-1) {stations};
\node[entidad] (mag) at (1,-1) {measurement\_types};
\node[entidad] (cat) at (2,-1) {sensor\_categories};
\node[entidad] (zon) at (3,-1) {zones};
\node[entidad] (cha) at (0,-2) {channels};
\node[entidad] (sen) at (1,-2) {sensors};
\node[entidad] (typ) at (2,-2) {sensor\_types};
\node[entidad] (loc) at (3,-2) {locations};
\node[entidad] (obs) at (0,-3) {observations};
\node[entidad] (lnk) at (1,-3) {channel\_sensors};
\draw[relacion] (ver) -- (src);
\draw[relacion] (sta) -- (src);
\draw[relacion] (sta.south) -- ++(0,-0.22) -| ([xshift=-3mm]loc.west) -- (loc.west);
\draw[relacion] (mag) -- (src);
\draw[relacion] (typ) -- (cat);
\draw[relacion] (sen) -- (typ);
\draw[relacion] (sen) -- (sta);
\draw[relacion] (sen.north) -- (1,-1.5) -- (2.5,-1.5)
 -- (2.5,-0.5) -- (2,-0.5) -- (org.south);
\draw[relacion] (cha) -- (sta);
\draw[relacion] (cha) -- (mag);
\draw[relacion] (lnk) -- (cha);
\draw[relacion] (lnk) -- (sen);
\draw[relacion] (obs) -- (cha);
\draw[relacion] (obs.east) -- ++(0.10,0) |- ([yshift=-4mm]ver.south) -- (ver.south);
\draw[relacion] (loc) -- (zon);
\draw[relacion] (loc.east) -- ++(0.15,0) -- ++(0,2.5) -| (ver.north);
\draw[relacion] (zon) edge[loop right] (zon);
\end{tikzpicture}
\caption{Modelo relacional completo de catorce tablas. Cada flecha parte de
la relación hija y apunta a la relación referenciada. Usuarios es independiente.}
\end{figure}
\clearpage

El E--R con atributos y cardinalidades se conserva en
\archivo{docs/arquitectura/diagrama_3nf.md}.

Las relaciones admiten cero o muchos hijos, salvo canal--asignación, limitada
a una fila por canal. Zona, distrito padre, versión de inventario, ubicación de
estación, tipo de dispositivo y organismo son referencias opcionales; las demás
FK exigen un padre por hijo. Una zona local puede pertenecer a un distrito.
\archivo{location_geography} reconstruye ambos nombres, sin repetirlos en cada
ubicación. Los 21 nombres de distrito publicados no acreditan vigencia histórica.

\subsection{Lectura del modelo y cardinalidades}
Una estación identifica un punto lógico de la fuente de datos. Un canal identifica
la serie de una magnitud en esa estación, y una observación registra su valor en
un instante. No se equipara un canal con un dispositivo físico: esa asociación
se almacena por separado en \archivo{channel_sensors}, únicamente cuando se conoce.
En la generación observada actual no hay asignaciones físicas acreditadas.

La tabla siguiente complementa las flechas del diagrama. Cada fila describe la
referencia desde la tabla hija hacia su padre. Una referencia obligatoria exige
exactamente un padre; una opcional permite cero o uno. Cada padre admite cero o
muchas filas hijas, salvo la asignación de canal indicada al pie.

\begin{longtable}{>{\raggedright\arraybackslash}p{0.25\textwidth}>{\raggedright\arraybackslash}p{0.25\textwidth}>{\raggedright\arraybackslash}p{0.25\textwidth}>{\raggedright\arraybackslash}p{0.14\textwidth}}
\toprule Tabla hija & Clave foránea & Tabla padre & Referencia \\
\midrule\endhead
\archivo{source_versions} & \archivo{source_id} & \archivo{sources} & Obligatoria \\
\archivo{zones} & \archivo{parent_zone_id} & \archivo{zones} & Opcional \\
\archivo{locations} & \archivo{zone_id} & \archivo{zones} & Opcional \\
\archivo{locations} & \archivo{inventory_version} & \archivo{source_versions} & Opcional \\
\archivo{stations} & \archivo{source_id} & \archivo{sources} & Obligatoria \\
\archivo{stations} & \archivo{location_id} & \archivo{locations} & Opcional \\
\archivo{measurement_types} & \archivo{source_id} & \archivo{sources} & Obligatoria \\
\archivo{sensor_types} & \archivo{category_id} & \archivo{sensor_categories} & Obligatoria \\
\archivo{sensors} & \archivo{sensor_type} & \archivo{sensor_types} & Opcional \\
\archivo{sensors} & \archivo{station_id} & \archivo{stations} & Obligatoria \\
\archivo{sensors} & \archivo{organization_id} & \archivo{organizations} & Opcional \\
\archivo{channels} & \archivo{station_id} & \archivo{stations} & Obligatoria \\
\archivo{channels} & \archivo{measurement_type} & \archivo{measurement_types} & Obligatoria \\
\archivo{channel_sensors} & \archivo{channel_id} & \archivo{channels} & Obligatoria \\
\archivo{channel_sensors} & \archivo{sensor_id} & \archivo{sensors} & Obligatoria \\
\archivo{observations} & \archivo{channel_id} & \archivo{channels} & Obligatoria \\
\archivo{observations} & \archivo{source_version} & \archivo{source_versions} & Obligatoria \\
\bottomrule
\end{longtable}

\archivo{channel_sensors.channel_id} es también clave primaria: un canal admite
cero o una asignación, mientras que un sensor puede estar asociado a varios
canales. \archivo{users} no tiene claves foráneas en este modelo; no se dibujan
relaciones con usuarios que el esquema no implementa. La autorreferencia de
\archivo{zones} permite un distrito padre opcional para una zona local; el trigger
comprueba el nivel del padre.

\subsection{Dependencias funcionales y tercera forma normal}
La condición de 3FN exige que, en cada DF no trivial \(X\to A\), \(X\) sea
superclave o \(A\) atributo primo. La cobertura de DF y las claves candidatas se
registran en \archivo{sprint1/postgres/dependencias_3nf.json}. Se calcularon los
cierres de todos los subconjuntos de atributos y se comparó el inventario con
las columnas reales de PostgreSQL. Con esa cobertura, las catorce relaciones
satisfacen 3FN y BCNF.

Estación--fuente y magnitud--fuente permanecen en sus catálogos; canales y
observaciones no repiten fuente. La asignación sensor--canal evita duplicar la
estación del dispositivo. Tipo--categoría se conserva en \archivo{sensor_types}.

La vista contractual calcula la hora local desde UTC, proyecta constantes y
obtiene calidad, estado temporal, rol e intervalo desde la política por fuente.
También repone el SHA del snapshot en la referencia original. La reconstrucción
conserva el DTO v2, sus identificadores y valores.

La reunión por claves de catálogo conserva la información; las asignaciones
opcionales usan LEFT JOIN. Los triggers comprueban fuentes, estación del sensor,
inmutabilidad y aceptación v2. Se acreditan ocho rechazos de operaciones, sin
aislar cada CHECK o FK; las pruebas positivas y PG07 utilizan rollback.

JSONB es un dominio SQL tipado indivisible; no se afirma normalización escalar
recursiva. Los NULL opcionales denotan ausencia y no se usan como claves. No se
infieren DF de coincidencias del lote o de condiciones solo para valores no nulos.

\subsection{Dispositivos, magnitudes y procedencia}
\archivo{sensor_types} representa tipos de dispositivos administrativos;
\archivo{measurement_types} representa variables y unidades medidas. No hay FK
entre magnitudes y categorías físicas. Los doce tipos demostrativos y sus
categorías son explícitamente sintéticos. Se conservan 5.133 sensores y usuarios
demo y un organismo demo, sin enlaces inventados con canales observados.

\subsection{Modelo documental y compatibilidad}
MongoDB conserva \archivo{readings}, \archivo{incidents}, \archivo{derived_events},
\archivo{alerts}, \archivo{organizations}, \archivo{users} y \archivo{sensors}.
Las lecturas reúnen valor, calidad, reloj y procedencia, con BSON Date e
\archivo{_id=measurement_id}. Los objetos \archivo{measurements} y
\archivo{provenance} representan información semiestructurada. Esa flexibilidad
y los patrones de acceso justifican el modelo documental. La recuperación de
información narrativa se demuestra por separado con documentos sintéticos, sin
atribuir contenido no estructurado a toda lectura por el solo hecho de usar BSON.

\archivo{observation_contract} conserva la interfaz compartida, identidad y
payload, sin cambiar MongoDB ni establecer FK o transacciones distribuidas entre
motores. Persisten los límites de sensores físicos, geografía y reglas demo.

\clearpage
\subsection{Diagrama documental de la generación Madrid v2}
\begin{figure}[ht]
\centering
\begin{tikzpicture}[x=1cm,y=1cm,
  doc/.style={draw,rounded corners,align=left,text width=4.55cm,
  inner sep=6pt,font=\footnotesize,anchor=north},
  ref/.style={-{Latex[length=2mm]},thick},
  lab/.style={font=\scriptsize,fill=white,inner sep=2pt}]
\node[doc] (r) at (0,0) {\textbf{readings} [V]\\
  \texttt{\_id = measurement\_id}: string\\
  \texttt{channel\_id, station\_id}: string\\
  \texttt{source\_id, source\_version}: string\\
  \texttt{observed\_at}: BSON Date\\
  \texttt{value}: número; \texttt{unit}: string/null\\
  \texttt{sensor\_id}: null; origen observado\\
  \texttt{raw\_reference, provenance}: objetos\\
  \texttt{measurements}: objeto opcional\\
  Calidad, zona y versión contractual.};
\node[doc] (e) at (5.4,0) {\textbf{derived\_events} [V]\\
  \texttt{\_id, measurement\_id}: string\\
  \texttt{rule\_version}: string\\
  \texttt{threshold, value}: números\\
  \texttt{unit}: string; \texttt{zone\_id}: string/null\\
  \texttt{observed\_at, generated\_at}: Date\\
  \texttt{data\_origin}: derived\\
  \texttt{official\_event}: false\\
  Derivación por regla de demostración.};
\node[doc] (a) at (10.8,0) {\textbf{alerts} [V]\\
  \texttt{\_id, event\_id}: string\\
  \texttt{zone\_id, rule\_version}: string\\
  \texttt{observed\_at, emitted\_at}: Date\\
  \texttt{data\_origin}: derived\\
  \texttt{replay}: true\\
  \texttt{delivery}: local\_collection\_only\\
  \texttt{official\_alert}: false\\
  No constituye notificación externa.};
\draw[ref] (e.west) -- (r.east);
\draw[ref] (a.west) -- (e.east);
\node[doc] (o) at (0,-7) {\textbf{organizations} [sin V]\\
  \texttt{\_id = organization\_id}: string\\
  \texttt{name}: string\\
  \texttt{responsible\_person\_or\_role}: string\\
  \texttt{evidence\_reference}: string\\
  \texttt{provenance}: objeto\\
  \texttt{data\_origin}: synthetic.};
\node[doc] (s) at (5.4,-7) {\textbf{sensors} [sin V]\\
  \texttt{\_id = sensor\_id}: string\\
  Nombre, categoría y estación: string\\
  \texttt{organization\_id}: string\\
  Instalación y fecha de estado: string\\
  \texttt{provenance}: objeto\\
  Origen sintético; sin \texttt{sensor\_type}.};
\node[doc] (u) at (10.8,-7) {\textbf{users} [sin V]\\
  \texttt{\_id = user\_id}: string\\
  \texttt{name, email}: string\\
  \texttt{consent\_reference}: string\\
  \texttt{demo\_only}: boolean\\
  \texttt{provenance}: objeto\\
  Sin vínculo usuario--sensor definido.};
\draw[ref] (s.west) -- (o.east);
\node[doc] (i) at (0,-12) {\textbf{incidents} [V]\\
  \texttt{\_id, category}: string\\
  \texttt{reported\_date}: string (2026)\\
  \texttt{zone\_id}: string/null\\
  Fuente y versión: string\\
  \texttt{raw\_reference}: objeto\\
  \texttt{person\_fields\_removed}: true.};
\node[align=left,text width=9.8cm,anchor=north west,font=\footnotesize] at (2.7,-12) {
  \textbf{Leyenda.} Las flechas representan referencias lógicas del contenido,
  no claves foráneas MongoDB. [V] indica validador existente; los tipos mostrados
  describen los documentos cargados, no una restricción de tipo sobre cada campo.
  Las restricciones efectivas se detallan en el cargador y el catálogo nativo.\\[4pt]
  Estación, canal, zona y categoría remiten a catálogos externos a estas siete
  colecciones. No hay vínculo entre sensores administrativos y lecturas públicas:
  \texttt{sensor\_id} es nulo. Los incidentes no se unen temporalmente a las lecturas.};
\end{tikzpicture}
\caption{Colecciones oficiales y referencias lógicas de Madrid v2.}
\end{figure}

\begin{longtable}{>{\raggedright\arraybackslash}p{0.23\textwidth}>{\raggedright\arraybackslash}p{0.67\textwidth}}
\toprule Colección & Validación e índices, además de \texttt{\_id} \\
\midrule\endhead
readings & strict/error; UNIQUE (channel\_id, observed\_at, timestamp\_role);
(source\_id, measurement\_type, observed\_at); (zone\_id, measurement\_type). \\
derived\_events & error; (rule\_version, observed\_at). \\
alerts & error; (zone\_id, emitted\_at). \\
incidents & error; (category). \\
Administrativas & Solo el índice \texttt{\_id}; sin validadores adicionales. \\
\bottomrule
\end{longtable}
\clearpage

